%% 
%% Copyright 2007-2020 Elsevier Ltd
%% 
%% This file is part of the 'Elsarticle Bundle'.
%% ---------------------------------------------
%% 
%% It may be distributed under the conditions of the LaTeX Project Public
%% License, either version 1.2 of this license or (at your option) any
%% later version.  The latest version of this license is in
%%    http://www.latex-project.org/lppl.txt
%% and version 1.2 or later is part of all distributions of LaTeX
%% version 1999/12/01 or later.
%% 
%% The list of all files belonging to the 'Elsarticle Bundle' is
%% given in the file `manifest.txt'.
%% 
%% Template article for Elsevier's document class `elsarticle'
%% with harvard style bibliographic references

%\documentclass[preprint,12pt,authoryear]{elsarticle}

%% Use the option review to obtain double line spacing
%% \documentclass[authoryear,preprint,review,12pt]{elsarticle}

%% Use the options 1p,twocolumn; 3p; 3p,twocolumn; 5p; or 5p,twocolumn
%% for a journal layout:
%% \documentclass[final,1p,times,authoryear]{elsarticle}
%% \documentclass[final,1p,times,twocolumn,authoryear]{elsarticle}
%% \documentclass[final,3p,times,authoryear]{elsarticle}
%% \documentclass[final,3p,times,twocolumn,authoryear]{elsarticle}
%% \documentclass[final,5p,times,authoryear]{elsarticle}
 \documentclass[final,5p,times,twocolumn,authoryear, draft]{elsarticle}

%% For including figures, graphicx.sty has been loaded in
%% elsarticle.cls. If you prefer to use the old commands
%% please give \usepackage{epsfig}

%% The amssymb package provides various useful mathematical symbols
\usepackage{amssymb}
\usepackage{lipsum}
\usepackage{gensymb}

\usepackage{bm}

%% The amsthm package provides extended theorem environments
%% \usepackage{amsthm}

%% The lineno packages adds line numbers. Start line numbering with
%% \begin{linenumbers}, end it with \end{linenumbers}. Or switch it on
%% for the whole article with \linenumbers.
%% \usepackage{lineno}

%% You might want to define your own abbreviated commands for common used terms, e.g.:
\newcommand{\kms}{km\,s$^{-1}$}
\newcommand{\msun}{$M_\odot}

\journal{High Energy Astrophysics}


\begin{document}

\begin{frontmatter}

%% Title, authors and addresses

%% use the tnoteref command within \title for footnotes;
%% use the tnotetext command for theassociated footnote;
%% use the fnref command within \author or \affiliation for footnotes;
%% use the fntext command for theassociated footnote;
%% use the corref command within \author for corresponding author footnotes;
%% use the cortext command for theassociated footnote;
%% use the ead command for the email address,
%% and the form \ead[url] for the home page:
%% \title{Title\tnoteref{label1}}
%% \tnotetext[label1]{}
%% \author{Name\corref{cor1}\fnref{label2}}
%% \ead{email address}
%% \ead[url]{home page}
%% \fntext[label2]{}
%% \cortext[cor1]{}
%% \affiliation{organization={},
%%            addressline={}, 
%%            city={},
%%            postcode={}, 
%%            state={},
%%            country={}}
%% \fntext[label3]{}

\title{Title of paper}

%% use optional labels to link authors explicitly to addresses:
%% \author[label1,label2]{}
%% \affiliation[label1]{organization={},
%%             addressline={},
%%             city={},
%%             postcode={},
%%             state={},
%%             country={}}
%%
%% \affiliation[label2]{organization={},
%%             addressline={},
%%             city={},
%%             postcode={},
%%             state={},
%%             country={}}

\author[first]{Author name}
\affiliation[first]{organization={University of the Moon},%Department and Organization
            addressline={}, 
            city={Earth},
            postcode={}, 
            state={},
            country={}}

\begin{abstract}
%% Text of abstract
Example abstract for the High Energy Astrophysics Journal. Here you provide a brief summary of the research and the results.
\end{abstract}

%%Graphical abstract
%\begin{graphicalabstract}
%\includegraphics{grabs}
%\end{graphicalabstract}

%%Research highlights
%\begin{highlights}
%\item Research highlight 1
%\item Research highlight 2
%\end{highlights}

\begin{keyword}
%% keywords here, in the form: keyword \sep keyword, up to a maximum of 6 keywords
keyword 1 \sep keyword 2 \sep keyword 3 \sep keyword 4

%% PACS codes here, in the form: \PACS code \sep code

%% MSC codes here, in the form: \MSC code \sep code
%% or \MSC[2008] code \sep code (2000 is the default)

\end{keyword}


\end{frontmatter}

%\tableofcontents

%% \linenumbers

%% main text

\section{Introduction}
...

\section{Theory}
...

\section{Simulation Setup}
Simulations were performed using the VPIC (Vector Particle-In-Cell) code from LANL (Los Alamos National Laboratory). The Simulations were in 3D geometry ($N_x \times N_y \times N_z$) but we set $y$ domain size to unity, e.g. $N_y = 1$. A shock wave is excited by the so-called \textit{injection method}. A plasma consisting of electrons and ions is injected from the right-hand side boundary with an initial velocity $\bm{v}_0$ and travels toward to the negative $x$-direction.
At the left-hand side boundary, particles are reflected by the conducting wall. Then, a shock forms due to the interactions between the incoming and the reflected particles, and it propagates in the x-direction with a shock speed $v_{sh}$. Therefore, the simulations are done in the downstream rest frame. The periodic
boundary conditions are imposed in the y and z-direction.


Parameter units are normalized to natural PIC units: charges are in electron charge unit $e=1$, masses are in electron mass unit $m_e=1$, velocities are in speed of light unit $c=1$, time is normalized by electron plasma frequency $\omega_{pe}^{-1} = 1$, length is normalized by electron inertial length $d_e = c \omega_{pe}^{-1}$ and permittivity of space $\epsilon_0 = 1$. 

The domain spans from $0$ to $L_x$ in $x$, $0$ to $L_y$ in $y$ and $0$ to $L_z$ in $z$. The resolution is $L_x / N_x = L_y / N_y = L_z / N_z = 0.1$. The mesh size (topology) is chosen such that the kinetic scale is well resolved. The number of particles per cell (nppc) is $100$. The ion-electron mass ratio $m_i / m_e = 100$. The ratio of the plasma frequency to the electron cyclotron
frequency is $\omega_{pe} / \Omega_{ce} = 5$. In the upstream frame, electron and ion have a temperature $T_0 = 0.01$ (in units of $m_e c^2$) and number density $n_0=1$. We use $B_0= m_e c \Omega_{ce} / e$ for the upstream magnetic field



We consider three simulation cases that vary the shock normal angle $\theta$ (relative to the ambient magnetic field) and the injection velocity $\bm{v}_0$ as follows: (i) $\bm{v}_0 = 0.08c$ with $\theta = 75^\circ$, (ii) $\bm{v}_0 = 0.16c$ with $\theta = 75^\circ$, and (iii) $\bm{v}_0 = 0.16c$ with $\theta = 90^\circ$.



\begin{table}[h!]
	\centering
	\caption{Simulation parameters.}
	\setlength{\tabcolsep}{3pt} % Tightens the gaps between columns
	\begin{tabular}{p{0.52\columnwidth} ll}
		\hline
		Description & Parameter & Value \\ \hline
		Domain x size    & $N_x$                       & $8192$          \\
		Domain y size    & $N_y$                       & $1$             \\
		Domain z size    & $N_z$                       & $1024$          \\
		Simulation x size & $L_x$                     & $819.2 d_e$     \\
		Simulation y size & $L_y$                     & $0.1 d_e$       \\
		Simulation z size & $L_z$                     & $102.4 d_e$     \\
		Simulation time     & $\tau$                      & $100 \Omega_{ci}^{-1}$ \\
		Number of Particles per cell        & nppc               & $100$           \\
		Ion-to-electron mass ratio & $m_i / m_e$                & $100$           \\
		Plasma-to-cyclotron ratio & $\omega_{pe} / \Omega_{ce}$ & $5$    \\
		Initial plasma temperature & $T_i, T_e$                 & $0.01 m_e c^2$  \\
		Initial number density    & $n_{0,i}, n_{0,e}$          & $1$             \\
		Upstream magnetic field   & $B_0$                       & $0.2$           \\ 
		Ion inertial length   & $d_i$                       & $0.1 d_e$           \\
		\hline
	\end{tabular}
\end{table}

TODO:
\begin{enumerate}
	\item - beta parameters
	\item - Alven speed
	\item - Alven Mach number
\end{enumerate}

\section{Result}
...

\section{Discussion}
...

%\section*{Acknowledgements}
%Thanks to ...
%
%%% The Appendices part is started with the command \appendix;
%%% appendix sections are then done as normal sections

\appendix

\section{Filtering functions.}
...

\end{document}

\endinput
